\section{Introduction}\label{sec:intro}
Outdoor walkers and runners may wish to understand their routes, activity and nearby points of interest (POIs) without carrying additional wearable sensors. A smartphone offers positioning and motion inputs, while geographic information system (GIS) services can supply maps, POIs and route data. A 2025 Taiwan survey reported exercise participation of 83.3\% and regular participation of 35.6\% \cite{survey}; this motivates research on exercise support but does not establish that a particular system improves adherence.

Using phone data presents challenges: positioning drift and missing samples distort distance and speed; the nearest POI in a straight line may not be reachable by the shortest route; and continuous sampling consumes battery and storage. Energy expenditure is an estimate dependent on the activity, time and personal inputs. Speed or motion data alone cannot establish heart rate or cardiac workload.

We propose a mobile web experience in which an AI avatar helps plan a route, the user confirms it, and the phone records the activity. The design also includes food queries and confirmed records through a Model Context Protocol (MCP) tool, consented cloud storage, and a separately authorized digest-anchoring interface. The contribution is a unified proposed workflow, a lightweight sensing and GIS formulation, and a plan for evaluating accuracy, resource use, interaction and integrity. These are research objectives, not demonstrated health benefits.

Section~\ref{sec:related} discusses related work; Section~\ref{sec:method} presents the proposed scheme; Section~\ref{sec:experiments} describes development status and planned evaluation; and Section~\ref{sec:conclusion} concludes.
